\documentclass[11pt, a4paper]{exam}
\usepackage[a4paper,top=2.3cm,bottom=2.8cm,left=2cm,right=2cm,headheight=15pt,headsep=0.6cm,footskip=1.2cm]{geometry}
\usepackage[T1]{fontenc}
\usepackage[utf8]{inputenc}
\usepackage{lmodern}
\usepackage{textcomp}
\usepackage[french]{babel}
\usepackage{xcolor}
\usepackage{listings}
\usepackage{graphicx}
\usepackage{array}
\usepackage{enumitem}
\usepackage{amsmath}
\usepackage{tikz}
\usetikzlibrary{positioning,shapes.multipart,arrows.meta}
\usepackage[most]{tcolorbox}

\definecolor{codegreen}{rgb}{0,0.6,0}
\definecolor{codegray}{rgb}{0.5,0.5,0.5}
\definecolor{codepurple}{rgb}{0.58,0,0.82}
\definecolor{backcolour}{rgb}{0.95,0.95,0.92}

\lstdefinestyle{mystyle}{
    backgroundcolor=\color{backcolour},
    commentstyle=\color{codegreen},
    keywordstyle=\color{magenta},
    numberstyle=\tiny\color{codegray},
    stringstyle=\color{codepurple},
    basicstyle=\ttfamily\footnotesize,
    breakatwhitespace=false,
    breaklines=true,
    postbreak=\mbox{\textcolor{codegray}{$\hookrightarrow$}\space},
    captionpos=b,
    keepspaces=true,
    columns=fullflexible,
    numbers=left,
    numbersep=5pt,
    showspaces=false,
    showstringspaces=false,
    showtabs=false,
    upquote=true,
    tabsize=4,
    xleftmargin=1.5em,
    framexleftmargin=0.3em,
    aboveskip=0.6em,
    belowskip=0.6em,
    extendedchars=true,
    literate={é}{{\'e}}1 {è}{{\`e}}1 {ê}{{\^e}}1 {ë}{{\"e}}1 {à}{{\`a}}1 {â}{{\^a}}1
             {ç}{{\c c}}1 {ù}{{\`u}}1 {û}{{\^u}}1 {ô}{{\^o}}1 {î}{{\^i}}1 {ï}{{\"i}}1
             {É}{{\'E}}1 {È}{{\`E}}1 {À}{{\`A}}1 {Ç}{{\c C}}1 {œ}{{\oe}}1
             {«}{{\guillemotleft}}1 {»}{{\guillemotright}}1 {→}{{$\rightarrow$}}1
}
\lstset{style=mystyle, language=Python}
\lstdefinestyle{sqlstyle}{style=mystyle, language=SQL, morekeywords={LIMIT,REFERENCES,IN,TEXT,FLOAT,DATE,VARCHAR,CHAR,REPLACE}, numbers=none, xleftmargin=0.3em, frame=lines}

% Code en ligne
\newcommand{\code}[1]{\texttt{#1}}

% Titre d'une question : numéro et intitulé, barème à droite
\newcommand{\Qt}[2]{\subsection*{#1\hfill{\normalfont\small #2}}}
% Rappel de la consigne en une ligne
\newcommand{\rappel}[1]{\noindent\textit{Ce qui est demandé :} #1\par\smallskip}

% La réponse attendue, mise en évidence
\newtcolorbox{reponse}[1][Réponse]{enhanced, breakable, colback=black!5, colframe=black, boxrule=0.9pt, arc=1mm,
  title={#1}, fonttitle=\bfseries\small, coltitle=black, colbacktitle=black!15, left=2mm, right=2mm, top=1mm, bottom=1mm,
  before skip=0.5em, after skip=0.7em}
% Erreur fréquente
\newtcolorbox{erreur}{enhanced, breakable, colback=white, colframe=black, boxrule=0.5pt, leftrule=3pt, arc=0mm,
  title={Erreur fréquente}, fonttitle=\bfseries\small, coltitle=black, colbacktitle=white, left=2mm, right=2mm, top=1mm, bottom=1mm,
  before skip=0.5em, after skip=0.7em}
% Encadré « À retenir »
\newtcolorbox{retenir}{enhanced, breakable, colback=white, colframe=black, boxrule=0.5pt, arc=1mm, boxsep=1mm,
  borderline={0.4pt}{1.2pt}{black, dashed},
  title={À retenir}, fonttitle=\bfseries\small, coltitle=black, colbacktitle=white, left=2mm, right=2mm, top=1mm, bottom=1mm,
  before skip=0.5em, after skip=0.7em}

\newcommand{\methode}{\noindent\textbf{Méthode.}\ }
\newcommand{\explication}{\noindent\textbf{Explication.}\ }

\setlist{itemsep=0.15em, topsep=0.3em}
\renewcommand{\arraystretch}{1.25}
\setlength{\parindent}{0pt}
\setlength{\parskip}{0.35em}

% Piles dessinées verticalement (sommet en haut) ; \suppr barre un élément supprimé
\newcommand{\suppr}[1]{\ooalign{#1\cr\hidewidth\rule[0.55ex]{1em}{1pt}\hidewidth\cr}}
\newcommand{\pilev}[2]{\begin{tikzpicture}[baseline=(titre.base)]
  \node[font=\scriptsize, align=center, text width=22mm] (titre) at (0,0) {#1};
  \foreach \v [count=\k] in {#2} {\node[draw, minimum width=8mm, minimum height=5mm, inner sep=0pt, font=\small] at (0,-0.25-\k*0.5) {\v};}
\end{tikzpicture}}
\setlength{\emergencystretch}{3em}
% Remarque sur l'énoncé (erreur ou ambiguïté du sujet)
\newtcolorbox{remarque}{enhanced, breakable, colback=white, colframe=black, boxrule=0.5pt, arc=0mm,
  borderline={0.4pt}{1.2pt}{black, dotted},
  title={Remarque sur l'énoncé}, fonttitle=\bfseries\small, coltitle=black, colbacktitle=black!8, left=2mm, right=2mm, top=1mm, bottom=1mm,
  before skip=0.5em, after skip=0.7em}

\title{Corrigé — CC1 Term NSI - Rattrapage : Pile, File et Rappels de 1er}
\author{Babash - NSI Tani Malandi et Johanne Bakalara}
\date{Rattrapage de septembre 2025 (45 minutes, sur 20 points)}

\begin{document}
\pagestyle{headandfoot}
\firstpageheader{}{}{}
\runningheader{\small Corrigé}{}{\small Rattrapage pile, file et rappels -- 2025}
\runningheadrule
\firstpagefooter{}{\footnotesize Babash - NSI Tani Malandi et Johanne Bakalara \textperiodcentered{} CC BY-SA 4.0 \textperiodcentered{} page \thepage}{}
\runningfooter{}{\footnotesize Babash - NSI Tani Malandi et Johanne Bakalara \textperiodcentered{} CC BY-SA 4.0 \textperiodcentered{} page \thepage}{}

\maketitle
\thispagestyle{headandfoot}
\begin{center}\footnotesize\itshape D'après l'exercice 1 du sujet 22-NSIJ1G11 (baccalauréat 2022), le sujet 22-NSIJ2ME1 (baccalauréat 2022, Métropole, jour 2) et l'exercice 1 du sujet 21-NSIJ2G11 (baccalauréat 2021).\end{center}

\begin{tcolorbox}[colback=white, colframe=black, boxrule=0.5pt, arc=1mm, left=2mm, right=2mm]
\small\textbf{Comment lire ce corrigé.} Pour chaque question : ce qui est demandé, puis la réponse attendue (encadrée), puis l'explication pas à pas. Essaie d'abord seul·e, puis compare ta démarche, pas seulement ton résultat. Tous les codes de ce corrigé ont été exécutés en Python 3.
Le barème est celui du sujet (7 questions, 20 points).
\end{tcolorbox}

\begin{remarque}
Les titres des exercices annoncent 4 points (exercice I) et 6 points (exercice II). Les points donnés question par question font en réalité 12 points pour l'exercice I et 8 points pour l'exercice II, soit 20 points en tout. Ce corrigé suit les points des questions, comme le tableau de barème du sujet.
\end{remarque}

\section*{Exercice I -- Dictionnaires et listes (12 points)}

\begin{lstlisting}
jours = ["dimanche", "lundi", "mardi", "mercredi",
         "jeudi", "vendredi", "samedi"]
mois = {1 :("janvier",31), 2 :("février",28), 3 :("mars",31),
        4 :("avril",30)}
\end{lstlisting}

\begin{retenir}
Dans une \textbf{liste}, chaque élément a un \textbf{indice} (sa position), et le premier indice est \textbf{0} :
\begin{center}\small
\begin{tabular}{|l|*{7}{c|}}\hline
indice & 0 & 1 & 2 & 3 & 4 & 5 & 6 \\\hline
élément & dimanche & lundi & mardi & mercredi & jeudi & vendredi & samedi \\\hline
\end{tabular}
\end{center}
Dans un \textbf{dictionnaire}, on accède à une valeur par sa \textbf{clé}. Ici, la clé est le numéro du mois et la valeur est un \textbf{tuple} (nom, nombre de jours) : \code{mois[3]} vaut \code{("mars", 31)}.
\end{retenir}

\Qt{Question 1 -- Manipulation de liste et de dictionnaire}{2 points}

\textbf{a)} \rappel{obtenir \code{"jeudi"} à partir de la liste \code{jours}. \hfill\textit{(1 point)}}
\begin{reponse}
\code{jours[4]}
\end{reponse}
\explication \code{"jeudi"} est le cinquième élément. On compte à partir de 0 : dimanche 0, lundi 1, mardi 2, mercredi 3, jeudi 4.

\begin{erreur}
Écrire \code{jours[5]} en comptant à partir de 1 : on obtient \code{"vendredi"}.
\end{erreur}

\textbf{b)} \rappel{obtenir le \textbf{nom} du mois numéro 2 à partir de \code{mois}. \hfill\textit{(1 point)}}
\begin{reponse}
\code{mois[2][0]} (cette expression vaut \code{"février"}).
\end{reponse}
\explication
\begin{enumerate}
\item \code{mois[2]} : la valeur associée à la clé 2, le tuple \code{("février", 28)} ;
\item \code{[0]} : le premier élément de ce tuple, le nom \code{"février"}. (\code{mois[2][1]} donnerait le nombre de jours, 28.)
\end{enumerate}

\Qt{Question 2 -- Questions à choix}{4 points}

\textbf{a)} \rappel{le principal avantage d'un dictionnaire sur une liste. \hfill\textit{(1 point)}}
\begin{reponse}
\textbf{B.} Le dictionnaire permet un accès direct par clé, sans parcourir les éléments.
\end{reponse}
\explication Avec \code{mois[3]}, Python trouve directement la valeur de la clé 3. Pour trouver « mars » dans une liste de mois, il faudrait la parcourir. A est faux : la liste n'est pas plus rapide pour retrouver un élément par son nom. C est faux : les \textbf{clés} d'un dictionnaire sont uniques (si on réutilise une clé, l'ancienne valeur est remplacée), alors qu'une liste accepte les doublons.

\textbf{b)} \rappel{que renvoie \code{jours[7]} ? \hfill\textit{(1 point)}}
\begin{reponse}
\textbf{C.} Une erreur : \code{IndexError: list index out of range}.
\end{reponse}
\explication La liste a 7 éléments, d'indices 0 à 6. L'indice 7 n'existe pas : Python s'arrête avec l'erreur « indice de liste hors limites ». Le dernier élément est \code{jours[6]} (ou \code{jours[-1]}).

\textbf{c)} \rappel{que signifie « une liste est mutable » ? \hfill\textit{(1 point)}}
\begin{reponse}
\textbf{C.} Ses éléments ou sa taille peuvent être modifiés après sa création.
\end{reponse}
\explication \og Mutable \fg{} veut dire « qui peut changer ». Exemples : \code{jours[0] = "Dimanche"} modifie un élément ; \code{jours.append("fin")} change la taille. Un tuple, lui, n'est pas mutable. A est faux : une liste peut contenir des textes, comme \code{jours}. B décrit le contraire de « mutable ».

\textbf{d)} \rappel{combien de couples clé/valeur après \code{mois[5] = ('mai', 31)} ? \hfill\textit{(1 point)}}
\begin{reponse}
\textbf{B.} 5 couples clé/valeur.
\end{reponse}
\explication La clé 5 n'existe pas encore : l'affectation \textbf{ajoute} le couple \code{5 : ('mai', 31)}. Le dictionnaire passe de 4 à 5 couples (\code{len(mois)} vaut 5). Un dictionnaire est mutable : on peut lui ajouter des couples après sa création, donc C est faux. Si la clé existait déjà, sa valeur serait remplacée et la taille ne changerait pas.

\Qt{Question 3 -- La fonction \code{nb\_jours\_mois}}{2 points}
\rappel{renvoyer le nombre de jours du mois \code{numero\_mois}.}
\begin{reponse}
\begin{lstlisting}
def nb_jours_mois(mois, numero_mois):
    return mois[numero_mois][1]
\end{lstlisting}
\end{reponse}
\explication
\begin{enumerate}
\item \code{mois[numero\_mois]} donne le tuple du mois, par exemple \code{("mars", 31)} pour 3.
\item L'élément d'indice 1 de ce tuple est le nombre de jours.
\item On le \textbf{renvoie} avec \code{return}.
\end{enumerate}
Exemple : \code{nb\_jours\_mois(mois, 3)} renvoie \code{31} ; \code{nb\_jours\_mois(mois, 2)} renvoie \code{28}.

\begin{erreur}
Écrire \code{print(mois[numero\_mois][1])} au lieu de \code{return} : la valeur s'affiche, mais la fonction renvoie \code{None}. On ne peut alors pas réutiliser le résultat dans un calcul.
\end{erreur}

\Qt{Question 4 -- Code à trous : \code{changement\_mois}}{4 points}
\rappel{compléter la fonction qui renvoie le tuple (mois suivant, année suivante).}
\begin{reponse}
\begin{lstlisting}
def changement_mois (numero_mois, annee):
    if numero_mois == 12:
        return (1, annee + 1)
    else:
        return (numero_mois + 1, annee)
\end{lstlisting}
\end{reponse}
\explication
\begin{itemize}
\item Ligne 2 : le seul cas particulier est décembre, le mois 12.
\item Ligne 3 : après décembre vient janvier (1), \textbf{et} l'année augmente de 1.
\item Ligne 5 : sinon, on passe au mois suivant et l'année ne change pas.
\end{itemize}
\begin{center}\small
\begin{tabular}{|l|c|l|}\hline
Appel & \code{numero\_mois == 12} ? & Valeur renvoyée \\\hline
\code{changement\_mois(1, 2025)} & non & \code{(2, 2025)} \\\hline
\code{changement\_mois(12, 1900)} & oui & \code{(1, 1901)} \\\hline
\end{tabular}
\end{center}

\begin{erreur}
Écrire \code{if numero\_mois = 12:} avec un seul \code{=}. Le signe \code{=} range une valeur dans une variable ; pour comparer, il faut \code{==}. Python affiche une \code{SyntaxError}.
\end{erreur}

\section*{Exercice II -- Pile et file (8 points)}

\Qt{Question 5 -- Question de cours : LIFO et FIFO}{3 points}

\textbf{a)} \rappel{que signifie FIFO ? Pile ou file ? \hfill\textit{(2 points)}}
\begin{reponse}
FIFO veut dire \textit{First In, First Out}, c'est-à-dire « premier entré, premier sorti ». Cela caractérise une \textbf{file}.
\end{reponse}
\explication Dans une file, on ajoute les éléments à un bout et on les retire à l'autre bout : le premier arrivé est le premier servi. À l'inverse, une pile est LIFO (\textit{Last In, First Out}, « dernier entré, premier sorti »).

\textbf{b)} \rappel{un exemple simple de file dans la vie courante. \hfill\textit{(1 point)}}
\begin{reponse}
La file d'attente à la cantine ou à un guichet : la première personne arrivée est la première servie. (Autres exemples : les voitures qui attendent la barge, les documents en attente d'impression.)
\end{reponse}

\Qt{Question 6 -- Réduction de la pile par exécution du code}{2 points}
\rappel{le contenu de la pile 5 4 5 4 2 1 (du haut vers le bas) après l'appel de \code{parcourir\_\allowbreak pile\_\allowbreak en\_\allowbreak reduisant(P)}.}
\begin{reponse}
Après l'appel, la pile \code{P} contient, du haut vers le bas : \textbf{5 5 4 2 1}.
\end{reponse}

\methode On exécute le code à la main.
\begin{enumerate}
\item \code{reduire\_triplet\_au\_sommet} regarde les trois éléments du haut \code{a}, \code{b}, \code{c}. Si \code{a} et \code{c} ont la même parité (tous deux pairs ou tous deux impairs), \code{b} est supprimé.
\item Puis la boucle retire l'élément du sommet de \code{p} et le pose sur \code{q} : le triplet suivant commence un cran plus bas.
\item Quand \code{p} a moins de 3 éléments, on vide \code{p} dans \code{q}, puis on remet tout \code{q} dans \code{p} : l'ordre d'origine est rétabli.
\end{enumerate}

Trace (« P » pour pair, « I » pour impair ; les piles sont écrites du haut vers le bas) :
\begin{center}\small
\begin{tabular}{|l|l|l|l|l|}\hline
Étape & \code{p} avant & Triplet \code{a b c} & Parités de \code{a} et \code{c} & \code{p} et \code{q} après \\\hline
tour 1 & 5 4 5 4 2 1 & 5 4 5 & I et I : 4 supprimé & \code{p} = 5 4 2 1 ; \code{q} = 5 \\
tour 2 & 5 4 2 1 & 5 4 2 & I et P : rien & \code{p} = 4 2 1 ; \code{q} = 5 5 \\
tour 3 & 4 2 1 & 4 2 1 & P et I : rien & \code{p} = 2 1 ; \code{q} = 4 5 5 \\\hline
\multicolumn{5}{|l|}{\code{p} n'a plus que 2 éléments : fin de la première boucle.} \\\hline
2\textsuperscript{e} boucle & 2 1 & & & \code{p} vide ; \code{q} = 1 2 4 5 5 \\
3\textsuperscript{e} boucle & vide & & & \code{p} = 5 5 4 2 1 ; \code{q} vide \\\hline
\end{tabular}
\end{center}
Au tour 1, après la réduction, \code{p} vaut 5 5 4 2 1 ; on pose ensuite le 5 du sommet sur \code{q}.
Le parcours ne supprime qu'un élément : le 4 qui était entre les deux 5.

\begin{erreur}
Supprimer aussi le 4 de 5 4 2 (5 et 2 n'ont pas la même parité) ou oublier qu'on avance d'un cran après chaque triplet. Il faut suivre le code ligne par ligne, et noter \code{p} et \code{q} à chaque tour.
\end{erreur}

\Qt{Question 7 -- La fonction \code{est\_gagnant}}{3 points}
\rappel{renvoyer \code{True} si la pile a exactement deux éléments, \code{False} sinon.}
\begin{reponse}
\begin{lstlisting}
def est_gagnant(pile_finale):
    return taille(pile_finale) == 2
\end{lstlisting}
\end{reponse}
\explication \code{taille(pile\_finale) == 2} est une comparaison : elle vaut déjà \code{True} ou \code{False}, on peut la renvoyer directement. Variante, tout aussi juste :
\begin{lstlisting}
def est_gagnant(pile_finale):
    if taille(pile_finale) == 2:
        return True
    else:
        return False
\end{lstlisting}
Exemples : pour la pile 4 0, la fonction renvoie \code{True} ; pour la pile 5 5 4 2 1 de la question 6, elle renvoie \code{False}.

\begin{erreur}
Renvoyer le texte \code{"True"} (entre guillemets) : c'est une chaîne de caractères, pas le booléen \code{True}. Autre erreur : \code{taille(pile\_finale) >= 2}, qui déclarerait gagnante une pile de 5 éléments.
\end{erreur}

\end{document}
